\documentclass[11pt]{article}

\usepackage[T1]{fontenc}
\usepackage[utf8]{inputenc}
\usepackage{lmodern}
\usepackage{amsmath,amssymb,amsthm,mathtools}
\usepackage[a4paper,margin=28mm]{geometry}
\usepackage{microtype}
\usepackage{xurl}
\usepackage[hidelinks]{hyperref}
\usepackage{enumitem}
\usepackage{booktabs}
\usepackage{array}

\newtheorem{theorem}{Theorem}[section]
\newtheorem{lemma}[theorem]{Lemma}
\newtheorem{proposition}[theorem]{Proposition}
\newtheorem{corollary}[theorem]{Corollary}
\newtheorem{definition}[theorem]{Definition}
\newtheorem{remark}[theorem]{Remark}

\newcommand{\N}{\mathbb N}
\newcommand{\WIN}{\ensuremath{\mathrm{WIN}}}
\newcommand{\LOSS}{\ensuremath{\mathrm{LOSS}}}
\newcommand{\DRAW}{\ensuremath{\mathrm{DRAW}}}
\newcommand{\odd}{\operatorname{odd}}
\newcommand{\src}{\operatorname{src}}
\newcommand{\ct}{\operatorname{ct}}
\newcommand{\asl}{\operatorname{asl}}
\newcommand{\Moves}{\operatorname{moves}}

\title{The Two-Player $3n\mathbin{\pm}1$ Game:\
A Binary Normal Form and a Well-Founded No-Draw Proof}
\author{Grisha Pochuev\\
  \small\href{mailto:n_854@mail.ru}{n\_854@mail.ru}\\
  \small\url{https://github.com/Grisha-Pochuev/3n-plus-minus-1-game}}
\date{10 August 2026}

\begin{document}
\maketitle

\begin{abstract}
Consider the impartial two-player game on positive odd integers in which a
move from $n>1$ chooses $3n-1$ or $3n+1$ and removes all powers of two; reaching
$1$ wins immediately. Arbitrary play need not terminate: for example
$5\to 7\to 5$ is a legal cycle. The question is whether optimal play nevertheless
has finite remoteness from every starting position. We conjugate the game to a
binary normal form with an expanding move $A$ and a strictly contracting move
$B$, introduce constant-tail coordinates and a three-free source map, and rank
marked proof configurations by retained source, finite proof-height tokens, and
a typed fixed-fibre normalizer. The key entry step is a factor-removal lemma
that either removes one factor of three at the same tail exponent or produces
finite outcome provenance. This prevents an unranked return from a factor-free
three/two frame to the preceding exponent-one state. A one-shot initialization
then allows the typed normalizer to be entered without treating a larger
temporary inner source as a source decrease. The resulting well-founded marked
normalization excludes all draw positions.
\end{abstract}

\tableofcontents

\section{The game and the theorem}

The current state is a positive odd integer $n$. For $n>1$ the player to move
chooses one of
\[
  3n-1,\qquad 3n+1,
\]
and divides the chosen value by the largest possible power of $2$, leaving an
odd integer. If the resulting odd integer is $1$, the player who made the move
wins immediately. Players alternate with perfect information.

A position is called \WIN\ for the player to move if at least one move leads to
\LOSS. It is called \LOSS\ if every legal move leads to \WIN. These two classes
are the inductively generated finite-remoteness classes. Any position not in
either class is called \DRAW.

\begin{theorem}[Main theorem]\label{thm:main}
Every positive odd starting position has finite remoteness. Equivalently, the
$3n\pm1$ game has no \DRAW\ positions, and optimal play reaches $1$ after
finitely many moves.
\end{theorem}

The theorem is not a statement about arbitrary play. There are legal cycles,
such as
\[
  5\longrightarrow 7\longrightarrow 5,
\]
under suitable choices by the players.

\section{Binary conjugation}

Write an odd state as $n=2m+1$ and define
\[
  F(m)=\left\lceil\frac{3m}{2}\right\rceil .
\]
For a nonnegative integer $x$, let $R(x)$ be the integer obtained by deleting
the maximal alternating suffix of the binary expansion of $x$. Thus
$R(1001_2)=10_2$ and $R(10101_2)=0$.

\begin{proposition}[First normal form]\label{prop:first-normal}
For $m>0$, the two children in $m$-coordinates are exactly
\[
  F(m),\qquad F(R(m)).
\]
\end{proposition}

\begin{proof}
One has
\[
  3(2m+1)-1=2(3m+1),\qquad
  3(2m+1)+1=2(3m+2).
\]
Exactly one of $3m+1$ and $3m+2$ is odd. In the other expression, each
additional factor of two removed corresponds to crossing one additional
alternating bit at the end of the binary expansion of $m$. The process stops
at the first repeated adjacent bit, which is precisely the deletion defining
$R(m)$. Returning to $m$-coordinates gives the displayed pair.
\end{proof}

Since both children lie in the image of $F$, we represent the state $F(q)$ by
$q$. The conjugated game has terminal state $0$ and moves
\[
  A(q)=\left\lceil\frac{3q}{2}\right\rceil,
  \qquad
  B(q)=R(A(q)).
\]

\begin{proposition}[Strict contraction]\label{prop:B-contracts}
For every $q>0$,
\[
  B(q)<q.
\]
\end{proposition}

\begin{proof}
Deleting a nonempty binary suffix gives
\[
  B(q)\le \left\lfloor\frac{A(q)}2\right\rfloor
       \le \left\lfloor\frac{3q+1}{4}\right\rfloor<q
\]
for $q\ge2$, while $B(1)=0$.
\end{proof}

Thus all unbounded difficulty comes from the expanding branch $A$ and from the
unbounded alternating suffix removed by $B$.

\section{Constant-tail coordinates and coefficient sources}

For a positive odd integer $a$, an exponent $r\ge1$, and a phase
$e\in\{0,1\}$, put
\[
  Q_r^e(a)=a2^r-e.
\]
This is a binary word whose final constant block has length $r$ and consists
of the bit $e$.

\begin{lemma}[Long-tail recurrence]\label{lem:long-tail}
For every positive odd $a$, phase $e$, and $r\ge3$,
\[
  A(Q_r^e(a))=Q_{r-1}^e(3a),\qquad
  B(Q_r^e(a))=Q_{r-2}^e(3a).
\]
Moreover
\[
  B(Q_1^e(a))=B(Q_2^e(a)).
\]
\end{lemma}

\begin{proof}
For $r\ge3$, the expanding child still ends in at least two equal bits; hence
its maximal alternating suffix consists only of the final bit. This gives both
long-tail identities. At the boundary, the binary words of $3a-e$ and $6a-e$
differ by appending one bit that extends the same alternating suffix, so the
same prefix remains after deletion.
\end{proof}

Define the embedded three-free coefficient
\[
  J(s)=2A(s)+1.
\]
Every positive odd coefficient $a$ has a unique factorization
\[
  a=3^kJ(s),\qquad k\ge0.
\]
Indeed, remove all factors of $3$ and call the remaining odd integer $c$.
Then $c\not\equiv0\pmod3$ and $(c-1)/2$ is $0$ or $2$ modulo $3$.
The image of
\[
 A(2t)=3t,\qquad A(2t+1)=3t+2
\]
is exactly the set of nonnegative integers not congruent to $1$ modulo $3$,
and $A$ is injective. Hence there is a unique $s$ with
$A(s)=(c-1)/2$, equivalently $c=J(s)$. Unique factorization supplies the
unique exponent $k$. The integer $s$ is called the coefficient source of
$a$ and is denoted $\src(a)$.

\begin{lemma}[Source boundary]\label{lem:source-boundary}
For every $s>0$ there is a phase
\[
  \alpha(s)=1-\left(\left\lfloor\frac{s}{2}\right\rfloor\bmod2\right)
\]
such that
\[
  \odd(3J(s)+1-2\alpha(s))=J(A(s))
\]
with $2$-adic valuation one, while the opposite phase selects $J(B(s))$ with
valuation at least two.
\end{lemma}

\begin{proof}
Put $u=A(s)$. Then $J(s)=2u+1$ and
\[
  3J(s)+1-2e=2(3u+2-e).
\]
The bracket is odd exactly when $e=1-(u\bmod2)$. For $s=2t$ or $2t+1$ one
has $u\bmod2=t\bmod2$, so this phase is exactly $e=\alpha(s)$. If $u$ is
even, then $e=1$ and $3u+1=J(u)$; if $u$ is odd, then $e=0$ and
$3u+2=J(u)$. Thus the valuation-one child is $J(A(s))$. The opposite sign is
the other original $3n\pm1$ child of the odd state $J(s)$; under the
conjugacy of Proposition~\ref{prop:first-normal} that child is $J(B(s))$.
Since it is not the valuation-one branch, its valuation is at least two.
\end{proof}

Multiplying a constant-tail coefficient by $3$ preserves its coefficient
source. Passing through a factor-free signed boundary with coefficient $J(s)$
therefore exposes exactly one ordinary $A/B$ move on the source.

For a positive state $q$, write $\kappa(q)$ for the odd coefficient in its
canonical constant-tail coordinates.

\begin{lemma}[Coefficient pullback]\label{lem:coeff-pullback}
For every positive odd $w$,
\[
  \kappa(3w)=J(R(w)),\qquad
  \kappa(3w-1)=J(R(2w-1)).
\]
\end{lemma}

\begin{proof}
Let $m\ge1$ be the length of the maximal alternating suffix of the odd word
$w$, and put $z=R(w)$. If $z>0$, maximality says that the leading bit of the
suffix equals the last bit of $z$. Because the suffix ends in $1$, its value is
\[
 t_m=\begin{cases}
 (2^m-1)/3,&m\text{ even},\\
 (2^{m+1}-1)/3,&m\text{ odd}.
 \end{cases}
\]
Thus $w=2^mz+t_m$. If $m$ is even, $z$ is even and
\[
  3w+1=2^m(3z+1)=2^mJ(z).
\]
If $m$ is odd, $z$ is odd and
\[
  3w+1=2^m(3z+2)=2^mJ(z).
\]
If $z=0$, the whole word is alternating; then $m$ is odd and
$3w+1=2^{m+1}$, whose odd part is $J(0)=1$. Hence in all cases
$\kappa(3w)=\odd(3w+1)=J(R(w))$.

For the second identity set $v=2w-1$, which is odd. Since
$3v+1=2(3w-1)$, the first identity applied to $v$ gives
\[
 \kappa(3w-1)=\odd(3w-1)=\odd(3v+1)=J(R(v))=J(R(2w-1)).
\]
\end{proof}


\section{Finite outcomes, witnesses, and ranks}

The terminal state $0$ is \LOSS\ with height $h(0)=0$. A finite \WIN\ state has
height
\[
  h(x)=1+\min\{h(y):y\text{ is a \LOSS\ child of }x\},
\]
and a finite \LOSS\ state has height
\[
  h(x)=1+\max\{h(y):y\text{ is a child of }x\}.
\]
A state whose outcome has been proved but which is being carried only to
justify a later incidence is called a \emph{routing witness}. A
\emph{retained proof token} is a finite state for which the proof has also
established the comparison needed by the rank: it either seeds an empty inner
token set at a one-shot seed transition, or replaces a previously retained
token by certified descendants of strictly smaller height. This distinction is
important: a newly discovered finite state is never silently appended to a
ranked multiset.

\begin{lemma}[Multiset token rank]\label{lem:token-rank}
For a finite multiset $M$ of proof heights define
\[
  \Phi(M)=\mathop{\#}_{h\in M}\omega^h,
\]
the natural ordinal sum of the monomials $\omega^h$. Replacing one token of
height $H$ by any finite multiset of certified descendants of heights strictly
below $H$ strictly decreases $\Phi$.
\end{lemma}

\begin{proof}
In Cantor normal form the coefficient of $\omega^H$ decreases by one, while
all inserted terms have lower exponent. No finite collection of lower terms
can compensate for the loss at exponent $H$.
\end{proof}

Every positive state $q$ has unique canonical constant-tail coordinates
$q=Q_r^e(a)$ with $a$ odd and $r\ge1$. If
$a=3^kJ(s)$, define the \emph{state source} $\rho(q)=s$.

\begin{lemma}[Universal state-source bound]\label{lem:source-bound}
For every $q>0$,
\[
  \rho(q)\le \frac{q-1}{6}.
\]
\end{lemma}

\begin{proof}
Write $q=Q_r^e(3^kJ(s))$. Since $r\ge1$ and $e\le1$,
$q\ge2J(s)-1$. Moreover
$J(s)=2\lceil3s/2\rceil+1\ge3s+1$. Hence
$q\ge6s+1$.
\end{proof}

\begin{lemma}[Boundary coefficient descent]\label{lem:boundary-coeff}
Let $a$ be positive and odd and
\[
 p=B(Q_1^e(a))=B(Q_2^e(a)).
\]
If $p>0$, then its canonical coefficient satisfies
\[
  \kappa(p)\le\frac{3a+1}{4}.
\]
In particular $\kappa(p)<a$ for $a\ge3$.
\end{lemma}

\begin{proof}
The common child is $p=R(3a-e)$. Deleting a nonempty suffix gives
$p\le\lfloor(3a-e)/2\rfloor$. For every positive integer $y$, the odd
coefficient in its canonical constant-tail coordinates is at most $(y+1)/2$.
Therefore
\[
 \kappa(p)\le\frac{p+1}{2}\le\frac{3a+1}{4}.
\]
The final inequality is strict relative to $a$ for $a\ge3$.
\end{proof}

\section{Side identities and typed boundary objects}

The finite routing system uses two elementary binary identities. We record
them here because they are the only places where an exceptional residue class
enters.

\begin{lemma}[Ordinary side relation]\label{lem:side}
Put $S(q)=B(q)$ and $T(q)=B(A(q))$. If
\[
  q\not\equiv1,3,12,14\pmod{16},
\]
then
\[
  T(q)\in\{A(S(q)),B(S(q))\}.
\]
If $q$ is in one of the four exceptional classes, then $A(q)$ is not
exceptional.
\end{lemma}

\begin{proof}
Let $x=A(q)$ and let $k$ be the length of the maximal alternating suffix of
$x$. Direct reduction of $A(q)$ modulo $8$ shows that $k\ge3$ is possible
exactly for $q\equiv1,3,12,14\pmod{16}$. Outside these classes $k=1$ or
$2$. Writing $x$ as $S(q)$ followed by that one- or two-bit alternating word,
the formulas
\[
  A(2u)=3u,\qquad A(2u+1)=3u+2
\]
show that $A(x)$ is obtained from $A(S(q))$ by appending the corresponding
one- or two-bit alternating word. Deleting its maximal alternating suffix
therefore leaves either $A(S(q))$ or $B(S(q))$. The four exceptional residue
classes map under $A$ to $2,5,10,13\pmod{16}$, none of which is exceptional.
\end{proof}

\begin{corollary}[No four consecutive wins]\label{cor:no-four-win}
For no $q$ can all four states
\[
  q,A(q),A^2(q),A^3(q)
\]
be \WIN.
\end{corollary}

\begin{proof}
If $q$ and $A(q)$ are both \WIN, then $B(q)$ must be \LOSS; the same holds at
the next two positions. If $q$ is ordinary, Lemma~\ref{lem:side} makes
$B(A(q))$ a child of the \LOSS\ state $B(q)$, impossible. Hence $q$ must be
exceptional. Applying the same argument to $A(q)$ forces $A(q)$ exceptional,
contrary to Lemma~\ref{lem:side}.
\end{proof}

\begin{lemma}[Alternating lift identity]\label{lem:alt-lift}
For every positive integer $z$ there are $m\ge1$ and
$\delta=1-(z\bmod2)$ such that
\[
  3z+1=Q_m^\delta(J(R(z))).
\]
If $R(z)>0$, one may take $m$ equal to the length of the maximal alternating
suffix of $z$; if $R(z)=0$, one takes one more than the binary length of $z$.
\end{lemma}

\begin{proof}
First suppose $r=R(z)>0$, let $k$ be the alternating-suffix length and put
$\epsilon=r\bmod2$. Then
\[
 z=2^k r+t_{k,\epsilon},\qquad
 t_{k,0}=\Big\lfloor\frac{2^k}{3}\Big\rfloor,
 \quad t_{k,1}=\Big\lfloor\frac{2^{k+1}}{3}\Big\rfloor.
\]
Since $J(r)=3r+1+\epsilon$, direct substitution gives
\[
 3t_{k,\epsilon}+1+\delta=2^k(1+\epsilon),
\]
the two parity cases being exactly the two remainders of a power of $2$
modulo $3$. Hence
\[
 3z+1+\delta=2^kJ(r),
\]
which is the required identity with $m=k$. If $R(z)=0$, the whole binary word
of $z$ is alternating and starts with $1$, so
$z=t_{k,1}$. The same calculation gives
$3z+1+\delta=2^{k+1}=2^{k+1}J(0)$.
\end{proof}

\begin{lemma}[Exceptional side attachment]\label{lem:exceptional-side}
Let $q$ be exceptional and suppose
\[
 q\text{ is \WIN},\qquad A(q)\text{ is \DRAW},\qquad
 \ell=B(q)\text{ is \LOSS}.
\]
Then for some $m\ge1$ and $\delta\in\{0,1\}$ the two children of $A(q)$ are
exactly
\[
  Q_m^\delta(J(\ell)),\qquad Q_{m+1}^\delta(J(\ell)),
\]
and this adjacent pair contains a continuing \DRAW.
\end{lemma}

\begin{proof}
Write $q=16t+c$ with $c\in\{1,3,12,14\}$. The four direct exceptional
formulas are
\[
\begin{array}{c|c|c|c}
c&\ell&B(A(q))&A^2(q)\\ \hline
1&R(6t)&18t+1&36t+3\\
3&R(6t+1)&18t+4&36t+8\\
12&R(6t+4)&18t+13&36t+27\\
14&R(6t+5)&18t+16&36t+32.
\end{array}
\]
In every row put respectively $z=6t,6t+1,6t+4,6t+5$. Then
$\ell=R(z)$ and $B(A(q))=3z+1$. By Lemma~\ref{lem:alt-lift},
\[
  B(A(q))=Q_m^\delta(J(\ell))
\]
for some $m\ge1$. The last column is $2B(A(q))+\delta$, so it equals
$Q_{m+1}^\delta(J(\ell))$. These are precisely the two children of the
\DRAW\ state $A(q)$; therefore neither is \LOSS\ and at least one is \DRAW.
\end{proof}

For a source $x$ and phase $e$, let $O(x,e)$ denote the following exact
\emph{DRAW obligation}:
\begin{enumerate}[label=(\alph*)]
\item $Q_1^e(J(x))$ is \DRAW; or
\item there is a \DRAW\ state whose two children are exactly
      $Q_2^e(J(x))$ and $Q_1^e(J(x))$.
\end{enumerate}
This definition introduces no outcome assumption beyond the displayed
\DRAW\ state.

\begin{lemma}[Hidden parent for a typed three/two frame]\label{lem:hidden-parent}
Let $a$ be positive, odd and not divisible by $3$, and suppose
\[
  a\equiv1+e\pmod3.
\]
Put $H=Q_3^e(a)$, $G=Q_2^e(a)$ and
\[
  K=\frac{2H+e-1}{3}.
\]
Then $K$ is an integer and
\[
  \Moves(K)=\{H,G\}.
\]
\end{lemma}

\begin{proof}
For $e=0$ the congruence gives $H\equiv2\pmod3$ and
$K=(2H-1)/3$; for $e=1$ it gives $H\equiv0\pmod3$ and $K=2H/3$.
These are exactly the two inverse branches of the injective map
$A(q)=\lceil3q/2\rceil$. The last three bits of $H$ are equal, so the
contracting child removes only the last bit and equals $G$.
\end{proof}

\section{Fixed-tail factor removal and entry typing}

The next two lemmas isolate the only raw entry phenomenon that is not already
an obligation or a factor frame.

\begin{lemma}[Fixed-tail factor-removal diamond]\label{lem:fixed-tail-factor}
Let $c$ be positive and odd, $e\in\{0,1\}$, and $r\ge2$. Put
\[
  X=Q_r^e(3c),\qquad G=Q_r^e(c),\qquad H=Q_{r+1}^e(c),\qquad Y=A(G).
\]
Then
\[
  \Moves(H)=\{X,Y\}.
\]
If $X$ is \DRAW, either $G$ is \DRAW, or the diamond determines a finite
routing witness together with its exact incidence to $G,H,X,Y$ and the other
child of $G$.
\end{lemma}

\begin{proof}
Since $r+1\ge3$, Lemma~\ref{lem:long-tail} gives $A(H)=X$. For $r\ge3$ it
also gives $B(H)=Q_{r-1}^e(3c)=A(G)$; for $r=2$ the boundary identity gives
the same equality. Thus $B(H)=Y$.

Because $X$ is \DRAW, $H$ is not \LOSS. If $H$ is \WIN, then $Y$ is a
\LOSS\ child of $H$. Assume $H$ is \DRAW. Then $Y$ is not \LOSS. If $G$ is
\LOSS, then $Y$ is \WIN. If $G$ is \WIN, either $Y$ itself is a \LOSS\
witness, or the other child of $G$ is a \LOSS\ witness. These are all outcome
orientations. In every finite branch the witness and the parent-child
incidence just displayed are known exactly; no token-rank comparison is
asserted at this stage.
\end{proof}

\begin{definition}[Entry controls]\label{def:entry-controls}
A \emph{boundary entry} is an obligation $O(x,e)$ or, more generally, a
factor-free adjacent frame
\[
 \{Q_{m+1}^e(J(x)),Q_m^e(J(x))\},\qquad m\ge1,
\]
known to contain a continuing \DRAW. Its remaining tail length is recorded as
$\tau$. The typed three/two frame satisfying Lemma~\ref{lem:hidden-parent} is
a distinguished $m=2$ boundary constructor.
A \emph{loss-sibling entry} is one of the finitely many local constructors
in which a continuing \DRAW\ is accompanied by a finite \WIN/\LOSS\ routing
witness with a proved parent/child incidence: the fixed-tail diamond of Lemma~\ref{lem:fixed-tail-factor}, the exponent-one predecessor diamond of Lemma~\ref{lem:exp-one-lift}, an ordinary side diamond of Lemma~\ref{lem:side}, the exceptional side attachment of Lemma~\ref{lem:exceptional-side}, or the hidden-parent frame of Lemma~\ref{lem:hidden-parent}. The constructor records its remaining
constant-tail/factor exponent as the local counter $\tau$.
\end{definition}

\begin{lemma}[Entry typing]\label{lem:entry-typing}
Every finite branch produced by Lemma~\ref{lem:fixed-tail-factor} is a
loss-sibling entry in the sense of Definition~\ref{def:entry-controls}; no
ranked multiset is changed in order to make this classification.
\end{lemma}

\begin{proof}
The proof of Lemma~\ref{lem:fixed-tail-factor} gives the finite witness and
its exact incidence to $G,H,X,Y$ and the other child of $G$. These are
precisely the data of the first loss-sibling constructor. Classification is
therefore immediate; subsequent tail/factor normalization belongs to the
typed routing relation and is ranked there by $\tau$ and, when available,
by descendant-token comparisons.
\end{proof}

\begin{lemma}[Exponent-one lift]\label{lem:exp-one-lift}
Let $a$ be positive and odd and $k\ge1$. If
$R_k=Q_1^e(3^ka)$ is \DRAW, then
\[
  P_{k-1}=Q_2^e(3^{k-1}a)
\]
is either \DRAW\ or \WIN; in the latter case its other child is a finite
\LOSS\ witness.
\end{lemma}

\begin{proof}
The boundary recurrence gives $A(P_{k-1})=R_k$. A state with a \DRAW\ child
cannot be \LOSS. If $P_{k-1}$ is \WIN, the \DRAW\ child cannot witness the
win, so its other child is \LOSS.
\end{proof}

\begin{proposition}[Arbitrary factorful exponent-one entry]\label{prop:factorful-entry}
Let $a=J(s)$ and $k>0$. If $Q_1^e(3^ka)$ is \DRAW, then after finitely many
outcome-compatible local steps one obtains either a typed entry of
Lemma~\ref{lem:entry-typing}, or the factor-free exponent-two state
\[
  Q_2^e(a)\text{ \DRAW}.
\]
\end{proposition}

\begin{proof}
Apply Lemma~\ref{lem:exp-one-lift}. Its finite branch is a loss-sibling entry by Definition~\ref{def:entry-controls}. Otherwise $Q_2^e(3^{k-1}a)$ is \DRAW. Apply
Lemma~\ref{lem:fixed-tail-factor} with $r=2$ repeatedly. Whenever its finite
branch occurs, Lemma~\ref{lem:entry-typing} supplies a typed entry. If it never
occurs, one factor of $3$ is removed at the same tail exponent at each step,
so after $k-1$ steps the displayed factor-free exponent-two state is reached.
\end{proof}

\begin{proposition}[Complete exponent-one/two boundary reduction]\label{prop:boundary-reduction}
Let $s\ge0$, $k\ge0$, $e\in\{0,1\}$, and let
\[
  Q_r^e(3^kJ(s)),\qquad r\in\{1,2\},
\]
be \DRAW. For $s>0$, after finitely many local steps at least one of the following occurs: a typed entry is obtained, the retained source strictly
decreases, or the factor-free exponent-two state $Q_2^e(J(s))$ is \DRAW.
For $s=0$ the corresponding statement is Lemma~\ref{lem:zero-source}.
\end{proposition}

\begin{proof}
Assume $s>0$. There are four cases. If $r=1,k=0$, the state is the first
alternative of the obligation $O(s,e)$ and is already a boundary entry. If
$r=1,k>0$, use Proposition~\ref{prop:factorful-entry}. If $r=2,k=0$, we are
at the displayed factor-free exponent-two state. Finally, if $r=2,k>0$, apply
Lemma~\ref{lem:fixed-tail-factor} repeatedly with $r=2$: a finite branch is a
loss-sibling entry by Definition~\ref{def:entry-controls}, while otherwise
one factor of $3$ is removed at the same exponent at each step and after $k$
steps the factor-free exponent-two state is reached. These cases are disjoint
and exhaustive.
\end{proof}

\section{Factor-free exponent-two base entry}

Let $x>0$, $a=J(x)$ and
\[
  P=Q_1^e(a),\quad U=Q_2^e(a),\quad
  b=B(P)=B(U),\quad F=A(U)=Q_1^e(3a),\quad g=1-e.
\]
Assume $U$ is \DRAW\ and $x$ is the least coefficient source of any \DRAW. The common child $b$ cannot be $0$, because $0$ is \LOSS\ and a \DRAW\ has no \LOSS\ child. By Lemma~\ref{lem:boundary-coeff}, $b$ has canonical coefficient strictly below $J(x)$; because $J$ is strictly increasing, its state source is below $x$. Hence $b$ is not \DRAW.
As a child of a \DRAW\ it is not \LOSS, so
\[
  b\text{ is \WIN},\qquad F\text{ is \DRAW}.
\]
\begin{lemma}[First-factor anchor identity]\label{lem:first-factor}
For the states above,
\begin{equation}\label{eq:first-factor}
  9J(x)+1-2e=2^vJ(b),\qquad v\ge1.
\end{equation}
\end{lemma}

\begin{proof}
Put $a=J(x)$. If $e=0$, then $P=2a$, $A(P)=3a$ and $b=R(3a)$.
Apply the first identity of Lemma~\ref{lem:coeff-pullback} to the odd integer
$w=3a$:
\[
 \odd(9a+1)=\kappa(9a)=J(R(3a))=J(b).
\]
If $e=1$, then $P=2a-1$, $A(P)=3a-1$ and $b=R(3a-1)$. The second identity
of Lemma~\ref{lem:coeff-pullback}, again with $w=3a$, gives
\[
 \odd(9a-1)=J(R(6a-1)).
\]
By the boundary identity of Lemma~\ref{lem:long-tail},
$R(6a-1)=R(3a-1)=b$. Hence in both phases the odd part of
$9a+1-2e$ is exactly $J(b)$, which proves \eqref{eq:first-factor}.
\end{proof}

\begin{proposition}[Base-entry attachment]\label{prop:base-entry}
Every factor-free exponent-two \DRAW\ at a globally minimal source has a
finite continuation which either strictly lowers the retained source or
produces a boundary/loss-sibling typed entry.
\end{proposition}

\begin{proof}
Suppose first that $e=1-\alpha(x)$. The first source-boundary numerator is
divisible by $4$, while
\[
9J(x)+1-2e=3(3J(x)+1-2e)-2(1-2e)\equiv2\pmod4.
\]
Thus $v=1$. Writing $u=A(x)$ gives $b=3u+1$. The two children of the \DRAW\
state $F$ are
\[
  T=Q_1^g(J(b)),\qquad C=R(T),
\]
and direct appended-bit deletion gives
$C\in\{A(b),B(b)\}$. If $T$ is \DRAW, this is the boundary entry
$O(b,g)$. If $C$ is \DRAW, the other child $\ell$ of the \WIN\ state $b$ is a
finite \LOSS\ witness. If $b$ is ordinary, this is the ordinary-side
loss-sibling constructor. If $b$ is exceptional and $C=A(b)$,
Lemma~\ref{lem:exceptional-side} gives an adjacent \DRAW\ frame over
$\ell$, hence a boundary/loss-sibling entry. It remains only the exceptional
orientation $C=B(b)$. In that orientation
$b\le(9x+5)/2$ and the alternating suffix of $A(b)$ has length at least three,
so
\[
 C=B(b)\le\frac{A(b)}8\le\frac{3b+1}{16}.
\]
Lemma~\ref{lem:source-bound} then yields
\[
 \rho(C)\le\frac{C-1}{6}<\frac{3b+1}{96}
 \le\frac{27x+17}{192}<x,
\]
a strict retained-source exit.

Now suppose $e=\alpha(x)$. Lemma~\ref{lem:source-boundary} gives
$3J(x)+1-2e=2J(A(x))$, hence $v\ge2$. If $v\ge4$, then
\[
16J(b)\le9J(x)+1\le27x+19,
\]
and $J(b)\ge3b+1$ implies
\[
  b\le\frac{9x+1}{16}<x.
\]
Thus the continuing \DRAW\ has smaller retained source. If $v=2$, the two
children of $F$ are exactly $Q_2^g(J(b))$ and $Q_1^g(J(b))$, which is the
second form of $O(b,g)$. If $v=3$, they are the typed three/two frame
$Q_3^g(J(b)),Q_2^g(J(b))$. Reducing
\eqref{eq:first-factor} modulo $3$ gives
\[
  J(b)\equiv1+g\pmod3,
\]
so Lemma~\ref{lem:hidden-parent} applies. The hidden parent is a boundary
entry; its finite side, if one is exposed, is a loss-sibling entry. Hence no
untyped base row remains.
\end{proof}

\section{The source-zero branch}

For $k\ge0$, $r\ge1$ and $e\in\{0,1\}$ put
\[
  S_{k,r}^e=Q_r^e(3^k).
\]
These are precisely the constant-tail states with coefficient source zero.

\begin{lemma}[Zero-source boundary normalization]\label{lem:zero-source}
If a source-zero \DRAW\ exists, then after finitely many outcome-compatible
moves one reaches a typed boundary or loss-sibling entry. The one-shot entry
bit is not consumed before this finite normalization is complete.
\end{lemma}

\begin{proof}
For $r\ge3$, Lemma~\ref{lem:long-tail} gives
\[
  \Moves(S_{k,r}^e)=\{S_{k+1,r-1}^e,S_{k+1,r-2}^e\}.
\]
Following a \DRAW\ child strictly lowers $r$, so a \DRAW\ with $r\in\{1,2\}$
is reached after finitely many moves. If $k=0$ already at that boundary, the
four possibilities are explicit:
\[
 Q_1^1(1)=1\ \WIN,\quad Q_2^1(1)=3\ \WIN,\quad
 Q_1^0(1)=2\ \LOSS,\quad Q_2^0(1)=4\ \LOSS.
\]
Hence a boundary \DRAW\ must have accumulated at least one factor of $3$.

At the boundary write, for some $n\ge1$,
\[
  3^n+1-2e=2^jJ(y).
\]
The exact $2$-adic valuation is
\[
\begin{array}{c|cc}
 &n\text{ odd}&n\text{ even}\\ \hline
 e=0&j=2&j=1\\
 e=1&j=1&j=2+v_2(n).
\end{array}
\]
For $j\ge2$ the raw and signed exits form the adjacent factor-free frame
\[
  Q_{j-1}^{1-e}(J(y)),\qquad Q_j^{1-e}(J(y)),
\]
containing a continuing \DRAW; by Definition~\ref{def:entry-controls} this is a boundary entry, with its finite tail length recorded in $\tau$. For $j=1$ one has
\[
  y=\frac{3^{n-1}-1}{2}.
\]
When $e=0$ the raw exit is $A(y)$; when $e=1$ it is $B(y)$. The raw and
signed states are the two exits of the same boundary \DRAW\ configuration.
If the signed exit is \DRAW, it is already the exponent-one form of an
obligation. If the raw exit is the continuing \DRAW\ and the signed exit is
finite, that finite state is necessarily \WIN; choosing a canonical \LOSS\
child of it supplies an exact finite routing witness, hence a loss-sibling
entry. If both exits are \DRAW, use the signed obligation. Thus every row
reaches a typed entry after a finite pre-entry computation; no impossible
``source below zero'' comparison is used.
\end{proof}

\begin{lemma}[$A$-selecting side diamond]\label{lem:a-side}
Let $x>0$, $e=\alpha(x)$, $w=Q_1^e(J(x))$, and $y=A(x)$. Then
\[
  B(w)\in\{A(y),B(y)\}.
\]
\end{lemma}

\begin{proof}
The source-boundary identity gives
$A(w)=Q_1^{1-e}(J(y))=4A(y)+1+e$. For $e=0$ the appended bits are $01$;
for $e=1$ they are $10$. If the first appended bit repeats the last bit of
$A(y)$, deleting the alternating suffix leaves $A(y)$; otherwise the deletion
continues through the maximal alternating suffix of $A(y)$ and leaves
$B(y)$.
\end{proof}

\begin{lemma}[Universal $B$-transfer diamond]\label{lem:b-diamond}
Let $x>0$, let $e=1-\alpha(x)$, put $a=J(x)$ and $g=1-e$, and factor
\[
  3a+1-2e=2^jJ(y),\qquad j\ge2,
\]
so $y=B(x)$. Define
\[
 P=Q_1^e(a),\quad U=Q_2^e(a),\quad
 D=Q_j^g(J(y)),\quad C=Q_{j-1}^g(J(y)),\quad F=A(U).
\]
Then $D=A(P)$ and $C=B(P)=B(U)$. If $X$ is the common child of $D,C$ and
$Y$ is the other child of $C$, then
\[
  B(F)=Y.
\]
Consequently an obligation in the $B$-selecting phase transfers a continuing
\DRAW\ to the adjacent frame $\{D,C\}$; when $j=2$ the transferred frame is
again an obligation in the opposite phase.
\end{lemma}

\begin{proof}
The formulas for $D$ and $C$ are the signed boundary formula and the
shared-child identity. The remaining equality is a direct substitution in the
three cases $j=2$, $j=3$, and $j\ge4$: the constant suffix of $D$ determines
its common child $X$, while the one-step longer factor state $F$ has the same
raw side $Y$. For the outcome assertion, a common child of a \DRAW\ cannot be
\LOSS. If the exponent-one member is \DRAW, the frame $\{D,C\}$ already
contains a \DRAW. In the parent form of an obligation, the only remaining
orientation has the exponent-two member \DRAW; if $C$ were \WIN, outcome
recursion through the displayed diamond would force $Y$ both \LOSS\ and a
child of the \DRAW\ state $F$, impossible. Hence $C$ is \DRAW. For $j=2$,
$D,C$ are exactly the exponent-two/one pair over $J(y)$.
\end{proof}

\section{The typed routing system}

We now define the marked relation used after entry. A configuration contains
an outer retained source $s$, an outer token multiset $M$, an entry bit
$\eta\in\{1,0\}$, and, after entry, inner data
\[
  (a,\zeta,N,v,\tau).
\]
Here $a$ is the retained inner source anchor; $N$ is the finite multiset of
inner proof-token heights; $\zeta\in\{1,0\}$ is a one-shot token-seed bit;
$v$ is one of the finite control states below; and $\tau\in\N$ is a local
tail/factor counter. Temporary arithmetic cursors are not rank anchors. An inner source is promoted to the ranked anchor $a$ only when every already retained token remains a valid witness for the new typed state; otherwise the numerical value remains a cursor until a preceding token decrease permits a reset.

The bit $\eta=1$ means that the current retained-source fibre has not yet
initialized the typed relation. The strict transition $\eta:1\to0$ may
install the initial finite outer-token multiset $M$, an arbitrary finite inner
routing source $a$, and any already proved routing witnesses. This is why
$\eta$ is placed \emph{before} the token and inner-source components of the
rank. After $\eta=0$, $M$ can change only by certified strict descendant
replacement. If entry contains no token suitable for ranking, the inner
seed bit is set to $\zeta=1$; the first later certified token may be installed
in $N$ by the strict transition $\zeta:1\to0$. Once $\zeta=0$, every change of
$N$ is a certified strict descendant replacement.

Define
\[
  \Psi(N)=\mathop{\#}_{h\in N}\omega^h.
\]
The full lexicographic rank is
\begin{equation}\label{eq:full-rank}
  \mathcal R=(s,\eta,\Phi(M),\zeta,\Psi(N),a,c(v),\tau).
\end{equation}
Thus an entry-bit decrease may seed $M$ and $a$; a later token decrease may
reset the still-later source/control data. All natural components use their
usual well-order and both bits are ordered $1>0$.

The finite control set $\mathcal V$ consists of the fifteen names listed in Table~\ref{tab:control}; the names are labels for semantic constructors, not additional mathematical assumptions.

\begin{lemma}[A-obligation routing]\label{lem:a-routing}
An $A$-selecting obligation $O(x,\alpha(x))$ has only the following
outcome-compatible exits: a strict inner-source decrease, a strict inner-token
replacement, a bounded canonical $A$-streak, or a factor fork carrying an
actual finite witness.
\end{lemma}

\begin{proof}
Let $e=\alpha(x)$, $g=1-e$, $P=Q_1^e(J(x))$, $U=Q_2^e(J(x))$ and
$b=B(P)=B(U)$. Lemma~\ref{lem:a-side} makes $b$ an ordinary child of $A(x)$ and gives
\[
 \rho(b)\le\frac{9x+1}{24}<x.
\]
Thus a \DRAW\ at $b$ is a strict inner-source exit. Suppose $b$ is \WIN. If
$P$ is \DRAW, its other child is \DRAW\ and is exactly the next obligation
$O(A(x),1-e)$. In the parent form of $O(x,e)$, the exponent-one state is
\WIN, the exponent-two state is \DRAW, and outcome recursion forces an
actual \LOSS\ sibling; the other child of the exponent-two state is an
adjacent factor frame carrying that witness. These are the stated rows.
\end{proof}

\begin{lemma}[Bounded canonical $A$-streak]\label{lem:a-streak}
A canonical $A$-continuation from an $A$-selecting obligation performs at
most four equal-rank side tests before reaching a $B$-selecting control or a
strict source/token exit.
\end{lemma}

\begin{proof}
During the streak the side states are consecutive points
$x_2,x_3,x_4,x_5$ on one $A$-ray. If all four were \WIN, this would contradict
Corollary~\ref{cor:no-four-win}. The first non-\WIN\ side state is either a
\DRAW, giving a strict source exit in the ordinary case or the adjacent frame of
Lemma~\ref{lem:exceptional-side} in the exceptional case, or the phase has already become
$B$-selecting. Thus no fifth equal-rank test exists.
\end{proof}

\begin{lemma}[$B$-selecting transfer]\label{lem:b-routing}
Every \texttt{b\_select} control carries, as part of its provenance, a retained
anchor $p$ and a current $B$-selecting cursor $y$ with
\[
  y\le\frac{3p+1}{2}.
\]
Let $j\ge2$ be the selected signed valuation and $z=B(y)$. If $j\ge3$, then
\[
  z\le\frac{27p+7}{48}<p,
\]
so the retained inner source strictly decreases. If $j=2$, the transferred
frame is again an exact obligation. Its first occurrence may use the one-shot
token seed; every later compatible valuation-two recycle either strictly
lowers the retained inner source or replaces the retained witness by a
certified lower descendant.
\end{lemma}

\begin{proof}
Lemma~\ref{lem:b-diamond} sends the $B$-selecting obligation to an adjacent frame over $z=B(y)$. When $j\ge3$, at least two additional suffix
bits are deleted after the mandatory first division, and the exact source
estimate is
\[
 z\le\frac{9y-1}{24}
   \le\frac{27p+7}{48}<p.
\]
For $j=2$ the transferred exponent-two/one pair is exactly the same
obligation type in the opposite phase. The first occurrence exposes the
finite ordinary-side witness and, if no inner token has yet been seeded,
uses $\zeta:1\to0$. In a later valuation-two recycle the common-side identity
places the new witness below the currently retained finite witness in its
proof tree; hence its height is strictly smaller. No unrelated token is
inserted after $\zeta=0$.
\end{proof}

\begin{lemma}[Boundary-entry normalization]\label{lem:boundary-normalize}
A \texttt{boundary\_entry} either makes a strict source/token transition or,
after finitely many equal-rank tail steps, reaches
\texttt{a\_obligation} or \texttt{factor\_fork}.
\end{lemma}

\begin{proof}
For an obligation there is nothing to normalize. Otherwise the entry is an
adjacent factor-free frame with finite lower exponent $m$. Follow a continuing
\DRAW\ through the long-tail recurrence while the lower exponent is at least
three. Each pure routing step lowers the recorded tail counter $\tau$. At the
first exponent-one/two boundary, the shared-child/source identities give an
obligation if the coefficient is factor-free and a factor fork if powers of
$3$ have accumulated. The hidden-parent congruence handles the typed
three/two constructor without introducing an additional row. Thus the macro is
finite; any earlier source or certified token exit is already strict.
\end{proof}

\begin{lemma}[Loss-sibling normalization]\label{lem:loss-sibling}
A \texttt{loss\_sibling\_entry} either makes a strict token/source transition
or reaches \texttt{factor\_fork} after finitely many equal-rank steps.
\end{lemma}

\begin{proof}
The entry records one of the finite constructors of
Definition~\ref{def:entry-controls} and the remaining tail/factor length
$\tau$. If the constructor is already an exponent-one/two or hidden-parent
frame, it is a factor fork. Otherwise, choose the continuing \DRAW\ branch in
the long-tail diamond. Lemma~\ref{lem:long-tail} reduces the remaining tail
length. Whenever the carried finite witness meets the continuing branch in an
ordinary side diamond, Lemma~\ref{lem:side} either keeps the same witness
incidence or replaces it by an actual child, which is a strict proof-height
replacement. An exceptional side can occur for only one step because its
$A$-successor is ordinary. Thus an equal-rank passage strictly decreases
$\tau$ and must reach the exponent-one/two factor boundary. No new witness is
introduced during this countdown.
\end{proof}

\begin{lemma}[Factor and marked-tail routing]\label{lem:factor-routing}
Every factor-fork, high-return, marked-tail, short-lift, and terminal control
has one of the following effects:
\begin{enumerate}[label=(\alph*)]
\item strict inner-source decrease;
\item strict replacement of an existing inner token by certified lower
      descendant token(s);
\item a control change listed in Table~\ref{tab:control} with the retained
      source and token multiset unchanged; or
\item a same-control recurrence with strictly smaller local counter $\tau$.
\end{enumerate}
\end{lemma}

\begin{proof}
The boundary and loss-sibling controls are covered by Lemmas~\ref{lem:boundary-normalize} and~\ref{lem:loss-sibling}.

The factor exponent-one row uses the shared exponent-one/two child. If the
selected finite source is ordinary, Lemma~\ref{lem:side} exposes a common
child of the retained \LOSS\ witness and hence a lower \WIN\ token; the
exceptional row is an exact constant-tail lift over the same \LOSS\ witness.
For factor exponent at least two, the reverse constant-tail parent either
removes a factor at exponent at least two or exposes a finite parent-child
witness; in the latter case its common child is a certified descendant.

A high return has valuation $5,6$, or at least $7$. In the first two rows the
explicit short-tail formulas replace the carried return token before the new
marked task is installed. For valuation at least $7$, the universal
three/two frame carries a \WIN/LOSS\ pair; the selected common child is a
certified descendant before the tail is restarted. Hence every high-return
edge is strict in $\Psi(N)$.

For a marked constant tail of length $D\ge3$, if its marked states are
$q$ (\WIN) and $y$ (\LOSS), the common child
\[
  r=B(q)=A(y)
\]
for $D\ge4$ (and the boundary common child for $D=3$) is a \WIN\ child of
$y$, so $h(r)<h(y)$. Thus the existing token $y$ is replaced before any new
inner task is installed. The two short rows $D=1,2$ are exact lifts over an
already carried token and merely change control. Long exact-lift recurrences
strictly lower their remaining tail/factor exponent, which is $\tau$.
Terminal exponent-two/three macros reconstruct a finite descendant token
before they restart any obligation. The strict token replacements used in these rows are collected explicitly in Table~\ref{tab:tokens}. These cases exhaust the factor/tail constructors because their only unbounded parameters are the valuation and tail intervals just split above.
\end{proof}

The equal-retained-rank control edges and a concrete control height $c(v)$ are
shown in Table~\ref{tab:control}. Any edge not listed there is strict in an
earlier component of \eqref{eq:full-rank}. Same-control tail/factor recurrences
strictly decrease $\tau$.

\begin{table}[ht]
\centering
\small
\begin{tabular}{@{}lll@{}}
\toprule
Control $v$ & $c(v)$ & Equal-rank successors \\
\midrule
\texttt{marked\_tail} & 9 & \texttt{short\_lift} \\
\texttt{boundary\_entry} & 8 & \texttt{a\_obligation}, \texttt{factor\_fork} \\
\texttt{short\_lift} & 8 & \texttt{a\_obligation} \\
\texttt{a\_obligation} & 7 & \texttt{a\_test\_4}, \texttt{factor\_fork} \\
\texttt{a\_test\_4} & 6 & \texttt{a\_test\_3}, \texttt{b\_select} \\
\texttt{a\_test\_3} & 5 & \texttt{a\_test\_2}, \texttt{b\_select} \\
\texttt{a\_test\_2} & 4 & \texttt{a\_test\_1}, \texttt{b\_select} \\
\texttt{a\_test\_1} & 3 & \texttt{b\_select} \\
\texttt{b\_select} & 2 & \texttt{b2\_first} \\
\texttt{loss\_sibling\_entry} & 2 & \texttt{factor\_fork} \\
\texttt{b2\_first} & 1 & \texttt{b2\_ready} \\
\texttt{factor\_fork} & 1 & \texttt{high\_return} \\
\texttt{b2\_ready} & 0 & -- \\
\texttt{high\_return} & 0 & -- \\
\texttt{terminal\_macro} & 0 & -- \\
\bottomrule
\end{tabular}
\caption{Equal-rank control graph. Every listed edge strictly lowers $c(v)$.
Strict source/token/bit edges may reset $c$ and $\tau$.}
\label{tab:control}
\end{table}

\begin{proposition}[Well-founded typed normalizer]\label{prop:typed-normalizer}
The typed routing relation is well-founded for every typed entry.
\end{proposition}

\begin{proof}
Consider the rank \eqref{eq:full-rank}. A retained-source decrease is strict in
the first component and may reset all later data. At fixed retained source,
$\eta:1\to0$ is strict and may seed $M$ and all later routing data. Once
$\eta=0$, every change of $M$ is a certified strict descendant replacement
and therefore decreases $\Phi(M)$ by Lemma~\ref{lem:token-rank}; it may reset
all later inner data.

At fixed earlier components, $\zeta:1\to0$ is a one-time strict inner-token
seed. Thereafter every change of $N$ strictly decreases $\Psi(N)$. With the
token components fixed, a promoted inner-source comparison strictly lowers
$a$. If all preceding components are unchanged, the semantic routing lemmas
above show that a control-changing equal-rank edge is one of the edges of
Table~\ref{tab:control}, hence strictly lowers $c(v)$; a same-control
tail/factor recurrence strictly lowers $\tau$. Thus every transition
strictly decreases one component while preserving all earlier components.
The lexicographic product of these well-orders is well-founded.
\end{proof}

\section{Proof of the main theorem}

\begin{proof}[Proof of Theorem~\ref{thm:main}]
Assume that a \DRAW\ exists and choose the least state source $s$ occurring
among \DRAW\ positions. Begin with $\eta=1$; no ranked token has yet been installed.

If $s=0$, Lemma~\ref{lem:zero-source} performs a finite pre-entry
normalization and reaches a boundary or loss-sibling typed entry. Consume
$\eta:1\to0$ exactly there and invoke Proposition~\ref{prop:typed-normalizer}.

Let $s>0$. Follow an actual \DRAW\ child through long constant tails until
the tail exponent is one or two. Proposition~\ref{prop:boundary-reduction}
then gives a typed entry, a strict retained-source exit, or the factor-free
exponent-two state $Q_2^e(J(s))$ \DRAW. The last case is handled by
Proposition~\ref{prop:base-entry}. A strict smaller-source exit contradicts
the choice of $s$; otherwise a typed entry is obtained and $\eta$ is
consumed. Thus all four boundary combinations $(r,k)=(1,0),(1,>0),(2,0)$ and
$(2,>0)$ are covered.

From then on every macro transition lies in the well-founded relation of
Proposition~\ref{prop:typed-normalizer}, unless an earlier outer component
strictly decreases, in which case the lexicographic rank decreases even more
strongly. Yet every \DRAW\ has no \LOSS\ child and at least one \DRAW\ child.
Therefore from every marked \DRAW\ configuration an outcome-compatible
continuing \DRAW\ can be chosen, yielding another finite macro transition.
Iteration would create an infinite descending chain in the rank
\eqref{eq:full-rank}, impossible. Hence no \DRAW\ exists in the conjugated
game.

It remains only to transfer the conclusion back to an arbitrary original
odd starting state. Write it as $n=2m+1$. If $m=0$, then $n=1$ is terminal.
For $m>0$, Proposition~\ref{prop:first-normal} says that its two original
children are $F(m)$ and $F(R(m))$. These are represented in the conjugated game
by the states $m$ and $R(m)$, respectively, and both have finite remoteness by
the no-\DRAW\ conclusion just proved. Hence the original state $n$ is itself
\WIN\ or \LOSS\ with finite height. Thus every positive odd starting position
has finite remoteness and optimal play reaches $1$.
\end{proof}

\section{Verification boundary and reproducibility}

The theorem above is a human mathematical claim. The accompanying files have
narrower roles.
\begin{itemize}[leftmargin=*]
\item \texttt{verify\_v3\_0.py} checks the long-tail identities, fixed-tail
      diamonds, ordinary side relation, base-entry valuations and inequalities,
      source-zero valuation table, hidden-parent identity, and the finite
      equal-rank control graph over large finite ranges.
\item \texttt{routing\_certificate\_v3\_0.json} records the control states,
      equal-rank edges, strict edge classes, and rank components used in
      Table~\ref{tab:control}. It is a structural certificate, not a substitute
      for the semantic proofs in this article.
\item The repository's Lean development checks selected definitions and
      general well-founded metatheory. It is not claimed to kernel-check the
      complete concrete theorem.
\end{itemize}

The finite computations support the arithmetic and the transcription of the
routing table; the infinite conclusion rests on the proved identities,
outcome recursion, and the well-founded rank.

\appendix
\section{Exceptional side formulas}\label{app:exceptional}

For completeness, the four exceptional classes of Lemma~\ref{lem:side} have
exact formulas. For $t\ge0$,
\[
\begin{array}{c|c|c}
q&B(q)&B(A(q))\\ \hline
16t+1  &R(6t)   &18t+1\\
16t+3  &R(6t+1) &18t+4\\
16t+12 &R(6t+4) &18t+13\\
16t+14 &R(6t+5) &18t+16.
\end{array}
\]
The displayed $B(A(q))$ value always has constant-tail length one. Hence an
exceptional side step cannot repeat immediately; its next $A$-source is
ordinary. These formulas also justify the exceptional branch used in
Proposition~\ref{prop:base-entry}.

\section{Routing inventory}\label{app:routing}

The equal-rank graph in Table~\ref{tab:control} is deliberately smaller than
the full outcome case split: source and token exits are omitted from the graph
because they are already strict in earlier components. The semantic partitions
are as follows.
\begin{center}
\small
\begin{tabular}{@{}p{.28\linewidth}p{.62\linewidth}@{}}
\toprule
Control & Exhaustive nonterminal partition \\
\midrule
\texttt{boundary\_entry} & source/token exit,
  $\texttt{a\_obligation}$, or $\texttt{factor\_fork}$ after finite tail countdown \\
\texttt{loss\_sibling\_entry} & source/token exit, or
  $\texttt{factor\_fork}$ after finite tail countdown \\
\texttt{a\_obligation} & source exit, token exit, canonical $A$ test, or
  retained-witness factor fork \\
\texttt{a\_test\_4,3,2,1} & source exit, token exit, next test, or
  $B$-selecting switch; the last test has no next-test row \\
\texttt{b\_select} & valuation $2$, or valuation at least $3$ (strict source) \\
\texttt{b2\_first} & strict source, or one-time witness installation \\
\texttt{b2\_ready} & strict source, or strict witness descendant \\
\texttt{factor\_fork} & exponent-one source/token, higher source/token, or
  adjacent high return \\
\texttt{high\_return} & valuations $5,6$ or at least $7$; each spends a token
  before marked-tail restart \\
\texttt{marked\_tail} & $D=1$, $D=2$, or $D\ge3$; long row replaces the
  marked \LOSS\ token by its common \WIN\ child \\
\texttt{short\_lift} & exact lift into an obligation \\
\texttt{terminal\_macro} & exponent-two/three terminal rows; each begins
  after a strict token replacement \\
\bottomrule
\end{tabular}
\end{center}
This inventory is exactly the case structure used in the proof of
Proposition~\ref{prop:typed-normalizer}; there is no implicit ``other'' row.

\begin{table}[ht]
\centering
\small
\begin{tabular}{@{}p{.23\linewidth}p{.29\linewidth}p{.37\linewidth}@{}}
\toprule
Routing row & Retained replacement & Certified comparison \\
\midrule
ordinary factor fork & $b\mapsto X$ & $X$ is a child of the canonical
  \LOSS\ child of $b$, hence $h(X)\le h(b)-2$ \\
exceptional factor fork & $b\mapsto \ell$ & $\ell$ is the actual \LOSS\
  child of $b$, hence $h(\ell)=h(b)-1$ \\
high return & $(u,c)\mapsto(p,b)\mapsto(q,y)$ & each new state is an actual
  descendant of the token it replaces; each replacement is strict \\
marked tail $D\ge3$ & $y\mapsto r$ & $r$ is a child of the retained
  \LOSS\ state $y$, so $h(r)\le h(y)-1$ \\
terminal exponent $2/3$ & $\ell\mapsto Z_2$ or $Z_3$ & the reconstructed
  terminal token is a proper proof-tree descendant of $\ell$ \\
zero-recycle exponent $3$ & $y\mapsto r$ & $r$ is a child of the retained
  \LOSS\ token $y$ \\
\bottomrule
\end{tabular}
\caption{Token replacements used by the strict routing rows. A branching
exceptional diamond may replace one token by several descendants; in that
case Lemma~\ref{lem:token-rank} applies to the resulting multiset.}
\label{tab:tokens}
\end{table}

\begin{thebibliography}{9}
\bibitem{Guy1986}
R. K. Guy, John Isbell's Game of Beanstalk and John Conway's Game of
Beans-Don't-Talk, \emph{Mathematics Magazine} 59(5) (1986), 259--269.

\bibitem{Guy1996}
R. K. Guy, Unsolved Problems in Combinatorial Games, Problem 42, in
\emph{Games of No Chance}, MSRI Publications 29, Cambridge University Press,
1996.

\bibitem{Althofer2024}
I. Alth\"ofer, M. Hartisch, T. Zipproth, Analysis of a Collatz Game and Other
Variants of the $3n+1$ Problem, \emph{Advances in Computer Games} (2024),
123--132.

\bibitem{AlthoferPage}
I. Alth\"ofer, Collatz Problems: The Two-Player $3n+-1$ Game,
\url{https://althofer.de/collatz-prizes.html}.

\bibitem{Repo}
G. Pochuev, Optimal $3n\pm1$ Game: Proof and Verification Repository,
\url{https://github.com/Grisha-Pochuev/3n-plus-minus-1-game}.
\end{thebibliography}

\end{document}
